\stfhIndex{I082}{Fuluoerlianyingshe}{弗洛尔链映射（Floer chain map）}
\section{\personFloer{}链映射的定义}
\label{sec:def-floer-chain-map}

设给定一族 $(\mathcal H,\mathfrak j)$，其中
$\mathcal H=\{H^s\}_{0\leq s\leq1}$、
$\mathfrak j=\{J^s\}_{0\leq s\leq1}$，以及拉长函数
$\rho:\mathbb R\to[0,1]$。对每个这样的对子 $(\mathcal H,\mathfrak j)$ 与
$\rho$，$(\mathcal H,\mathfrak j;\rho)$ 携带\emph{非自治}方程 (18.5.42)，其渐近条件为
\begin{equation}
 u(-\infty)=z_0,\qquad u(\infty)=z_1,
\label{eq:19-chain-map-asymptotics}
\end{equation}
同伦条件为固定的 $[u]=C\in\pi_2(z_0,z_1)$。以
\[
 \mathcal M((\mathcal H,\mathfrak j;\rho);z_0,z_1;C)
\]
表示满足这些渐近边界条件与同伦条件的方程 (18.5.44) 的有限能量解集。

链同态
\[
 h_{\mathcal H}=h_{(\mathcal H,\mathfrak j;\rho)}:
 CF_*(H^0)\longrightarrow CF_*(H^1)
\]
通过遍历周期轨道对子 $z_0,z_1$ 与同伦类 $C$，考察这些模空间
$\mathcal M((\mathcal H,\mathfrak j;\rho);z_0,z_1;C)$ 来定义。容易调整定义 19.2.2
中的\personConley{}--\personZehnder{}指标定义，以得到链映射所需的指标。

\begin{exercise}
\label{ex:cz-index-continuation-triple}
调整定义 19.2.2，为连接 $H^0$ 与 $H^1$ 的同伦 $\mathcal H$ 所对应的三元组
$(z_0,z_1;C)$ 定义\personConley{}--\personZehnder{}指标，其中
$z_0\in\operatorname{Per}(H^0)$、$z_1\in\operatorname{Per}(H^1)$ 且
$C\in\pi_2(z_0,z_1)$。
\footnote{译者注：原书把 $z_1$ 与 $\operatorname{Per}(H^1)$ 的关系印作子集符号
$z_1\subset\operatorname{Per}(H^1)$；因 $z_1$ 是单条周期轨道，译文改为成员符号。}
\end{exercise}

\begin{definition}[链映射]
\label{def:floer-regular-continuation-data}
若下列条件成立，则称 $(\mathcal H^\rho,\mathfrak j^\rho)$ \emph{\personFloer{}正则}。
\begin{enumerate}[label=(\arabic*)]
\item 若 $\mu_{\mathcal H^\rho}(z_0,z_1;C)\leq-1$，则
\[
 \mathcal M(\mathcal H^\rho,\mathfrak j^\rho;z_0,z_1;C)=\varnothing.
\]
\item 对任意满足
\[
 \mu_{\mathcal H^\rho}(z_0,z_1;C)=0
\]
的 $z_0\in\operatorname{Per}(H^0)$ 和 $z_1\in\operatorname{Per}(H^1)$，
$\mathcal M(\mathcal H^\rho,\mathfrak j^\rho;z_0,z_1;C)$ 横截且紧致，因而为有限集。
记
\[
 n(\mathcal H^\rho,\mathfrak j^\rho;z_0,z_1;C)
 :=\#\mathcal M(\mathcal H^\rho,\mathfrak j^\rho;z_0,z_1;C)
\]
为其中元素的代数计数；其他情形令
$n(\mathcal H^\rho,\mathfrak j^\rho;z_0,z_1;C)=0$。
\item 对任意满足
\[
 \mu_{\mathcal H^\rho}(z_0,z_2;C)=1
\]
的 $z_0\in\operatorname{Per}(H^0)$ 和 $z_2\in\operatorname{Per}(H^1)$，
$\mathcal M(\mathcal H^\rho,\mathfrak j^\rho;z_0,z_2;C)$ 横截，并可紧化为光滑
一维流形，其边界由断裂流线 $u_1\#_\infty u_2$ 的全体组成；其中或者
\[
\begin{gathered}
 (u_1,u_2)\in
 \mathcal M(\mathcal H^\rho,\mathfrak j^\rho;z_0,y;C_1)
 \times\mathcal M(H^1,J^1;y,z_2;C_2),\\
 \mu_{\mathcal H^\rho}(z_0,y;C_1)=0,\qquad
 \mu_H(y,z_2;C_2)=1,
\end{gathered}
\]
或者
\[
\begin{gathered}
 (u_1,u_2)\in
 \mathcal M(H^0,J^0;z_0,y;C_1)
 \times\mathcal M(\mathcal H^\rho,\mathfrak j^\rho;y,z_2;C_2),\\
 \mu_H(z_0,y;C_1)=1,\qquad
 \mu_{\mathcal H^\rho}(y,z_2;C_2)=0,
\end{gathered}
\]
并且前一种断裂中 $y\in\operatorname{Per}(H^1)$，后一种断裂中
$y\in\operatorname{Per}(H^0)$；对所有相应的
$C_1\in\pi_2(z_0,y)$、$C_2\in\pi_2(y,z_2)$，有 $C_1\#C_2=C$
\footnote{译者注：原书本定义记号多处漂移：第 (1)--(2) 项把标题中的
$(\mathcal H^\rho,\mathfrak j^\rho)$ 写成 $(\mathcal H_{\mathbb R},\mathfrak j_{\mathbb R})$（下标为黑板体 $\mathbb R$）；第 (3) 项
又写未拉长的 $(H,J)$，并在第一类断裂的后一因子中把同伦类写成上标 $C^2$、
在第二类断裂的后一因子及同伦类中把 $C_2$ 误印为 $C_1$；第二类断裂的指标行
还把 $\mathcal H^\rho$ 与 $H$ 的槽位互错（印作
$\mu_{\mathcal H^\rho}(z_0,y;C_1)=1$、$\mu_H(y,z_2;C_2)=0$）。源书又通篇以
$\operatorname{Per}(H_0)$、$\operatorname{Per}(H_1)$ 与
$\operatorname{Crit}\mathcal A_{H_0}$、$\operatorname{Crit}\mathcal A_{H_1}$
的下标形式书写 $H^0,H^1$，并写 $y\in\operatorname{Per}(H)$，而同伦中没有单个
$H$。译文依本定义标题、紧邻的指标式、$C_1\#C_2=C$ 以及第一类断裂自身的
配对模式，统一为上标记号并配以相容的指标。}
。
\end{enumerate}
\end{definition}

对每个给定的
$[z_0,w_0]\in\operatorname{Crit}\mathcal A_{H^0}$ 与
$[z_1,w_1]\in\operatorname{Crit}\mathcal A_{H^1}$，定义
\[
\begin{aligned}
 &\mathcal M((\mathcal H,\mathfrak j;\rho);[z_0,w_0],[z_1,w_1])\\
 &\qquad:=\bigcup_C
 \mathcal M((\mathcal H,\mathfrak j;\rho);z_0,z_1;C),
\end{aligned}
\]
其中 $C\in\pi_2(z_0,z_1)$ 满足
\begin{equation}
 [z_1,w_1]=[z_1,w_0\#C].
\label{eq:19-chain-map-class-condition}
\end{equation}
这与~\eqref{eq:19-boundary-class-condition} 类似。若 $\rho$-拉长
$(\mathcal H^\rho,\mathfrak j^\rho)$ 在定义 19.3.2 的意义下\emph{\personFloer{}正则}，则称
$(\mathcal H,\mathfrak j;\rho)$ \emph{\personFloer{}正则}。

在定义 19.3.2 的条件下，可像边界映射那样，由矩阵元
$n_{(\mathcal H,\mathfrak j;\rho)}([z_0,w_0],[z_1,w_1])$ 定义次数为 $0$ 的映射
\[
 h_{(\mathcal H,\mathfrak j;\rho)}:CF(H^0)\longrightarrow CF(H^1).
\]
定义 19.3.2 的条件还蕴含 $h_{(\mathcal H,\mathfrak j)}$ 次数为 $0$，且满足
\[
 h_{(\mathcal H,\mathfrak j;\rho)}\circ\partial_{(H^0,J^0)}
 =\partial_{(H^1,J^1)}\circ h_{(\mathcal H,\mathfrak j;\rho)}.
\]
任意两个这样的链映射 $h_{(\mathfrak j_1,\mathcal H^1)}$ 与
$h_{(\mathfrak j_2,\mathcal H^2)}$ 也链同伦；下面给出证明。
